\chapter*{謝辞}
\addcontentsline{toc}{chapter}{\protect\numberline {} 謝辞}

% 本研究を進めるにあたり，お世話になった方々へ感謝の意を述べさせていただきます．
本研究を進めるにあたり，多くの方々に多大なる御支援と御指導を賜りました．この場をお借りして，心より感謝申し上げます．


まずは，多くの御指導，御鞭撻を賜わりました梶克彦教授に感謝致します．
目指すべき方向から逸れてしまった際に適切にご指摘・ご修正いただく一方で，自ら考え研究を進める機会をいただけたことを大変ありがたく思います．
かなりスケジュールに対する進捗が遅れていたにも関わらず，辛抱強く導いてくださり心より感謝申し上げます．
また，本研究の価値を確信を持って肯定いただけたことによりこうして自身の研究に自信を持って取り組むことができるようになりました．
ここに深く御礼申し上げます．

% 先行研究への感謝
次に先行研究である滞在ウォッチプロジェクトを推進してこられた皆様にお礼を述べさせてください．
本プロジェクトは私が所属研究室に興味を持つきっかけになったプロジェクトであり，実際に同システムがあったことで研究室への来訪頻度の向上に繋がったと考えております．
研究に取り組む前に滞在ウォッチの軽微な改良を行わせていただいたのですがそうした改良に携われるような柔軟な設計をされており，とても貴重な機会をいただけましてありがたく思います．
その中でも特に本プロジェクトに所属する戸川浩汰さん，林航平さんに感謝を致します．
お二人は本プロジェクトメンバのみでの議論の機会を設けてくださり，何かちょっとでもわからないことがあった際に気軽に質問できるような関係性を築いてくださりました．
戸川さんには研究を取り組む上でのスケジュール管理や論文の論理展開など研究者として参考になる部分が多くあり，身近にそのような存在がいてくださったことに深く感謝しております．
林さんには同氏の研究から本研究の着想を得たことに加え，先行研究者として多くの助言をいただいたほか，本稿の一部画像の作成にもご協力いただきました．
また，井上翔人さんを含めたお三方には本稿の概論にもご助言いただき，より良い状態での提出ができましたことを感謝申し上げます．
林さん，井上さんを中心とした本プロジェクトの今後の更なる発展を心よりお祈りしております．


また，本システムの完成にあたり多田隆人さんにもご助力いただきました．
実装時のミスをご指摘いただき本研究を進めることができたこと，深く感謝申し上げます．
本研究室のウェブサイトの更新にも大きく貢献いただき，心より感謝申し上げます．


さらに，石井美帆さんには本研究を進める前から大変お世話になりました．
共に参加した外部イベントにおいて実装した内容は，本システム実装時に大いに役立ちました．
そのほか，オープンキャンパスでの展示の際など積極的に議論いただき，自身の成長につながった部分がかなりありました．
また，水谷祐生さんと長尾夏希さんを含めたお三方と行った輪読会を通じて屋内測位への知見を深めることができました．
重ねて御礼申し上げます．


加えて，日頃から熱心に討論，助言してくださり，実証実験やアンケートにもご協力くださいました梶研究室のみなさんに深く感謝致します．
その中でも同期の皆様に触発され，研究に対するモチベーションを高めることができました．
特に本研究でも引用させていただいた中島瑞希さんと粟生将之さんに感謝申し上げます．
お二人が早くから論文発表に向けて真摯に取り組まれてきた姿は，自身の研究姿勢を省みる大きなきっかけとなりました．
また，昨年中島さんと牧野遥斗さん，森真弓さんと共に取り組んだスマートフォンアプリケーションの開発経験は，本研究の進め方の指針になりました．
感謝致します．
そのほか，研究室内でイベントを企画し，楽しく過ごすことに助力してくださった林さん，石井さん，勝井翔英さん，花田光駿さんにも深く感謝致します．


さらに，本稿に対し助言してくださいました皆様により一層の感謝を申し上げます．
特に加藤風真さんには，提出直前まで本稿の確認にお付き合いいただき，心より感謝しております．
また，本研究室に所属する皆様に加えて，WiNF 2025(Workshop on Informatics 2025)でご質問ならびにご指摘いただいた皆様にも感謝申し上げます．
本ワークショップでの発表の機会は改めて自身の研究を俯瞰して見つめ直すきっかけとなるとともに，ポスターに情報を整理・集約するための貴重な経験となりました．



最後に，日々の生活を支え，彩ってくださった家族，友人，恋人に最大の感謝を捧げます．
母には，日々の家事や食事の準備を通して生活を支えてもらい，常に安心して研究に取り組める環境を整えてくれたことに深く感謝しています．
父には，金銭的な支援に加え，研究に行き詰まった際には客観的な視点から助言をしてくれるなど多方面にわたり支えてもらいました．
また，弟には，何かあった際に話を聞いてもらい，精神的な支えとなってくれたことに感謝しています．
友人には，大学に通うにあたって講義へのモチベーションを高めていただいたり，日常のささやかな出来事や嬉しい報告を共有してもらうことで，自身も前向きな気持ちで過ごすことができました．
また，多方面で努力を重ねている友人の姿に刺激を受け，自身の研究や学業に対する意欲を高めることができました．
恋人との時間は，研究活動を継続する上での良い気分転換となり，多様な経験を得る貴重な機会となりました．

以上，多くの方々の支えにより本研究を成し遂げることができました．
改めて深く感謝申し上げます．

% WiNFへの感謝

% Local Variables: 
% mode: latex
% TeX-master: "root"
% End: 

% 別場面でお世話になったとして同一人物が複数回記載される場合がありますがご了承ください

% 梶先生
% 研究の方向性の修正をいただいたこと，自身の研究の価値を認めてくれたこと，進捗が遅い自分を辛抱強く導いていただいたことなどなど研究者の先輩としての偉大さ

% 滞在ウォッチプロジェクトの先駆者
% 所属研究室に興味を持つきっかけ，実際に活用して行くこともあった点，一部改造に携われる柔軟な設計を行なっていた点

% togawaさん
% プロジェクトミーティングで大変お世話に，本人のお茶目さに対しての研究者としての尊敬，学祭での出店時にお誘いいただいたこと，先行研究の恩，アブストコメント感謝

% hayashiさん
% プロジェクトミーティングで大変お世話に，研究室内でのインフラ関係で気軽に質問させてもらうことが多かった，一部画像の作成にも協力いただいて頭が上がらない，先行研究の恩

% hirotoくん
% アブストでもコメントいただいた，プロジェクトミーティングもだし行動変容の方向でこのプロジェクトを絶やさずにいてもらえたのが嬉しい

% tadaくん
% システムの実装時のミスの修正を手伝ってもらった，ウェブ

% ishiiちゃん
% HackU，LT会

% mizutaniさん

% natsukiくん

% makinoくん

% mizukichan
% アブストコメントもくれてた，感謝

% mayumichan

% nishiさん

% fumaさん
% アブストコメントもくれてた，最近料理し出してるのもきゅうりきっかけ

% hanadaくん

% 同期全体

% WiNF

% ほのちゃん
% サークルメンツ（TeamAI＋漫研）
% 北坂先生も？
% いずる
% 愛知組＋ネ友
% バイト先
% 恋人
% 両親（支えの母と客観視＆送迎の父）
% 弟